%! Exponential Function Context and Data Modeling — Unit 2, Lesson 2.5 — Homework
\documentclass[10pt]{article}
\usepackage{precalculus-article}
\usepackage{precalculus-boxes}
\newcommand{\TallMath}[1]{$\displaystyle #1\rule[-1.4em]{0pt}{3.2em}$}

\begin{document}

\pageheader{Unit 2, Lesson 2.5}{Homework}
\namedateperiod

\vspace{0.12in}

\begin{enumerate}[itemsep=18pt]

\item A substance decays at a rate of 8\% per year. Its initial amount is 300 grams.

      \begin{enumerate}[label=(\alph*), itemsep=6pt]
        \item Write a model $f(t) = ab^t$ for the amount remaining after $t$ years.

              \vspace{0.06in}
              $a = $ \blank{3cm} \quad $b = $ \blank{3cm} \quad $f(t) = $ \blank{5cm}

        \item Is this growth or decay? How do you know from the value of $b$? \blank{6cm}
      \end{enumerate}

\item An exponential function passes through $(2,\,20)$ and $(5,\,540)$. Find $a$ and $b$.
      Show all algebraic steps.

      \vspace{0.08in}
      Equation 1: \blank{5cm} \quad Equation 2: \blank{5cm}

      \vspace{0.06in}
      Divide to eliminate $a$: \blank{9cm}

      \vspace{0.06in}
      $b = $ \blank{3cm} \quad then $a = $ \blank{4cm}

      \vspace{0.04in}
      Model: $f(x) = $ \blank{7cm}

\item A radioactive sample has a half-life of 10 years. Its initial mass is 80 grams.

      \begin{enumerate}[label=(\alph*), itemsep=8pt]
        \item Write the exponential model $M(t)$ for the mass remaining after $t$ years.

              \vspace{0.06in}
              $M(t) = $ \blank{7cm}

        \item Predict the mass after 25 years. Show your work.

              \vspace{0.06in}
              $M(25) = $ \blank{8cm} $\approx$ \blank{2cm} grams

        \item Is $t = -5$ a meaningful input in this context? Explain.

              \vspace{0.06in}
              \writeline
      \end{enumerate}

\item The function $f(d) = 5 \cdot 3^d$ models a quantity that grows daily, where $d$ is the
      number of days.

      \begin{enumerate}[label=(\alph*), itemsep=8pt]
        \item Rewrite $f$ as a function of weeks $w$ (where $d = 7w$).

              \vspace{0.06in}
              $g(w) = $ \blank{8cm}

        \item What does the base of $g(w)$ tell you about the weekly growth of the quantity?

              \vspace{0.04in}
              \writeline

        \item Rewrite $f$ as a function of 2-day periods $n$ (where $d = 2n$).

              \vspace{0.06in}
              $h(n) = $ \blank{8cm} \quad Base: \blank{2cm}
      \end{enumerate}

\item \textbf{(AP-Style Multiple Choice)} A population is modeled by $f(t) = 500 e^{0.04t}$,
      where $t$ is in years. Which of the following best describes the approximate annual growth rate?

      \medskip
      \begin{enumerate}[label=(\Alph*), itemsep=4pt]
        \item 0.04\%
        \item 4\%
        \item 4.08\%
        \item 40\%
      \end{enumerate}

      \vspace{0.06in}
      Answer: \blank{1.5cm} \quad Reasoning: \writeline

\item \textbf{(Extension)} The table below shows temperature readings (in ${}^\circ$C) at time
      $t$ minutes in a cooling experiment.

      \begin{center}
      \renewcommand{\arraystretch}{1.5}
      \begin{tabular}{c|c}
      $t$ (min) & Temperature (${}^\circ$C) \\ \hline
      0  & 80 \\
      5  & 55 \\
      10 & 42.5 \\
      15 & 36.25 \\
      \end{tabular}
      \end{center}

      The temperature is approaching a room temperature of 25${}^\circ$C. Let $y = T - 25$
      (subtract room temperature).

      \begin{enumerate}[label=(\alph*), itemsep=6pt]
        \item Compute the adjusted values $y = T - 25$ and check whether they form an
              exponential pattern (compute successive ratios).

              \vspace{0.06in}
              \renewcommand{\arraystretch}{1.6}
              \begin{tabular}{c|c|c}
              $t$ & $y = T - 25$ & Ratio (successive) \\ \hline
              0  & \blank{2.5cm} & --- \\
              5  & \blank{2.5cm} & \blank{2.5cm} \\
              10 & \blank{2.5cm} & \blank{2.5cm} \\
              15 & \blank{2.5cm} & \blank{2.5cm} \\
              \end{tabular}

        \item Write an exponential model for $y(t)$, then write the model for $T(t)$.

              \vspace{0.06in}
              $y(t) = $ \blank{5cm} \quad $T(t) = $ \blank{7cm}
      \end{enumerate}

\end{enumerate}

\end{document}
